\section{Finite initialization and quadratic improvement}\label{sec:initialization}
We now initialize the boundary recurrence and show that it continues
indefinitely. The sixth boundary moment starts a sequence of residual
estimates at $(34,b_0)$, $(78,b_1)$, and $(142,b_2)$. A physical-purity
estimate at $(240,b_2)$ then starts the coupled update. The inputs are
four variational energies and thermal moments on at most six sites.

\begin{proposition}[Finite boundary data]\label{in:finite}
Let $\alpha=.4042$ and $P_*(x)=x^3-.14x^2-.18x$. For the
opposite-field Hamiltonian \eqref{tr:Hopp},
\begin{equation}\label{in:trials}
 E_0(H_m^{\rm opp})+am<\alpha\qquad(m=34,78,142,240).
\end{equation}
At inverse temperature $b_0$,
\begin{equation}\label{in:thermal}
 Y_6(b_0)<.027,\qquad c_0P_*(\lambda_0)^2<.003745.
\end{equation}
\end{proposition}
\begin{proof}
Proposition~\ref{apptr:bound} proves \eqref{in:trials} with explicit
integer trial vectors. Proposition~\ref{appth:bounds} proves
\eqref{in:thermal} by finite thermal recurrences with certified
rounding errors. The appendices give all matrix and boundary data and
the precise finite inequalities.
\end{proof}

\subsection{Starting residual and physical-purity bounds}
The thermal estimates bound the residual and leading contributions to
$Y_m(b_0)$, while the variational estimate supplies a lower bound for
one ground weight.
Keep $\ell,U$ from the periodic seed and define, for each of the four
lengths in \eqref{in:trials},
\begin{equation}\label{in:epsR}
 \epsilon(m)=.027e^{b_0\alpha}U^m\frac{.874^{m-6}}{\ell^6},
 \qquad R(m)=e^{.132}U^m+\epsilon(m)-1.
\end{equation}

\begin{lemma}\label{in:seedbounds}
At these four lengths,
\begin{equation}\label{in:seedineq}
 e_m(0)\le\epsilon(m),\qquad
 A(m,j)\ge\bigl(1+R(m)^{2^j}\bigr)^{-2}\quad(j\ge0).
\end{equation}
\end{lemma}
\begin{proof}
By \eqref{be:residual}, $\rho_0\le.874$, and
$\ell<\lambda_0\le U$,
\[
 \lambda_0^mD_m(0)
 \le U^m .874^{m-6}\frac{Y_6(b_0)}{\ell^6}
 <e^{-b_0\alpha}\epsilon(m).
\]
The trial bound supplies a physical ground weight greater than
$e^{-b_0\alpha}$ and hence $Y_m(b_0)>e^{-b_0\alpha}$.
Dividing gives the first inequality.

To bound the full physical partition function from above, factor
$P_*(x)=x(x-1/2)(x+9/25)$. On $[\ell,U]$ it is positive and
$xP_*'(x)<4P_*(x)$. For example the latter is the positivity of
$4P_*(x)-xP_*'(x)=x(x^2-.28x-.54)$ on this interval.
It follows that $x^m/P_*(x)^2$ is increasing for $m>8$.
Consequently \eqref{in:thermal} gives
\[
 e^{b_0\alpha}\lambda_0^mc_0
 \le U^m\frac{.003745e^{b_0\alpha}}{P_*(U)^2}
 <e^{.132}U^m.
\]
The last scalar comparison has left coefficient below $1.133226$,
whereas $e^{.132}>1.141108$. Combining the leading and residual
estimates yields $e^{b_0\alpha}Y_m(b_0)\le1+R(m)$.

Select one physical ground eigenvector; uniqueness is not needed.
Its unnormalized weight $g$ is greater than $e^{-b_0\alpha}$, so
the sum of the remaining weights divided by $g$ is at most $R(m)$.
In particular $R(m)\ge0$. At $b_j=2^jb_0$, the corresponding
rest-to-ground ratio is at most $R(m)^{2^j}$, since a sum of
nonnegative $2^j$th powers is at most the $2^j$th power of the sum.
The purity is at least the square of the normalized weight of this
chosen ground eigenvector. This proves the second inequality.
\end{proof}

In particular $\epsilon(34)<.0928<1$, proving $c_0>0$.
The estimates $\rho_1\le.953$ and $\rho_2\le.971$ now allow
two applications of Lemma~\ref{be:recurrence}:
\begin{align}
 e_{78}(1)&\le F\bigl(.953^{44},.093,(1+R(78))^{-2}\bigr)<.133,
                                                     \label{in:78}\\
 e_{142}(2)&\le F\bigl(.971^{64},.133,(1+R(142)^2)^{-2}\bigr)<.14.
                                                     \label{in:142}
\end{align}
The bound at $m=240$, $j=2$ gives $1-A(240,2)<.21$.
Thus the induction of Proposition~\ref{be:coupled} starts with
\begin{equation}\label{in:initial}
 E_2=.14,\qquad\eta_2=.21.
\end{equation}
Both new leading overlaps have also been proved positive by the strict
defect bounds; no overlap was assumed at the next temperature.

\subsection{Explicit finite comparisons}
We record numerical margins to make the initialization reproducible.
The three lower bounds from Lemma~\ref{in:seedbounds} used above are
\begin{equation}\label{in:puritynumbers}
 \begin{split}
 (1+R(78))^{-2}&>.6180,\\
 (1+R(142)^2)^{-2}&>.7533,\\
 (1+R(240)^4)^{-2}&>.7923.
 \end{split}
\end{equation}
Using the full expressions in \eqref{in:epsR} gives upper outputs below
$.132129$, $.137301$, and $.207518$ in \eqref{in:78}, \eqref{in:142},
and $1-A(240,2)$ respectively. The lower bounds in
\eqref{in:puritynumbers} are rounded for display.

Starting from \eqref{in:initial}, the following rounded bounds satisfy
Proposition~\ref{be:coupled}:
\begin{equation}\label{in:table}
\begin{array}{c|cc|c}
 j&E_j&\eta_j&\text{upper bound on the $q$ used for step $j\to j+1$}\\\hline
 2&.14&.21&.142\\
 3&.12&.155&.0282\\
 4&.016&.026&.0028\\
 5&.001&.001&\text{---}
\end{array}
\end{equation}
The exponents $(M_j-k_j)/M_j$ are $98/240$, $1/2$, $1/2$.
The successive upper outputs for $E_{j+1}$ are below
$.117754,.015328,.000170$; those for $\eta_{j+1}$ are below
$.147078,.025188,.000432$. In computing the latter outputs one uses
the displayed rounded $E_{j+1}$, so each line closes with the next.

For clarity, all the scalar comparisons here and in the periodic
initialization have an elementary rational verification. For
$0\le t\le10$, put $T_{80}(t)=\sum_{r=0}^{80}t^r/r!$. Then
\begin{equation}\label{in:exp-enclosure}
 T_{80}(t)\le e^t\le T_{80}(t)+2t^{81}/81!.
\end{equation}
Indeed the ratio of consecutive omitted terms is at most $10/82<1/2$.
For negative exponents use reciprocal intervals. Rational positive
root comparisons can be checked by raising both sides to their
integer powers; for example the first $q$ bound in \eqref{in:table}
is $d(u_3)^{49}<.142^{120}$. Every denominator is positive.
These rules, exact rational input values, and the displayed formulas
suffice to verify the finite comparisons without floating-point
assumptions. The supplementary scalar certificate records all of them.

\subsection{Indefinite continuation}
We now enlarge the three defect bounds at $j=5$ to $.01$ so that a
single quadratic recurrence controls all later levels.

\begin{proposition}\label{in:quadratic}
Define $x_5=.01$ and $x_{j+1}=10x_j^2$ for $j\ge5$. Then, for every
$j\ge5$,
\begin{equation}\label{in:common}
 u_j\le x_j,\qquad e_{k_j}(j)\le x_j,\qquad
 1-A(M_j,j)\le x_j.
\end{equation}
In particular
\begin{equation}\label{in:xexplicit}
 x_j=10^{-1-2^{j-5}}.
\end{equation}
\end{proposition}
\begin{proof}
The table proves the two boundary inequalities at $j=5$, and
$u_5<.000015<.01$. Suppose \eqref{in:common} holds, and write
$x=x_j\le.01$. Here $k_j=M_j/2$ and $u_{j+1}=6u_j^2\le6x^2$.
The first argument of $F$ in \eqref{be:Eupdate} satisfies
\[
 \sqrt{d(u_{j+1})}\le
 \sqrt{\frac{6x^2}{2-18x^2}}\le2x.
\]
The first output is therefore bounded by
$F(2x,x,1-x)$. For $x>0$, division by $x^2$ gives
\[
 \frac{F(2x,x,1-x)}{x^2}
 =\frac{2}{1-2x}\bigl((1-x)^{-1}+(1-x)^{-2}+(1-x)^{-3}\bigr).
\]
This expression is increasing on $[0,.01]$ after its continuous
extension at zero. Its endpoint value is
$42430000/6792093<10$. Thus we may take $E_{j+1}=10x^2$.

For factors in $[0,1]$, the defect of their product is at most the sum
of their defects. Thus the four-factor product in \eqref{be:etaupdate}
has defect at most $x+x+6x^2+10x^2=2x+16x^2$. Moreover
\[
 \frac{f(2x+16x^2)}{x^2}
 =\frac{(2+16x)^2}{2(1-2x-16x^2)^2}
\]
is increasing on $[0,.01]$ and has endpoint value
$3645000/1495729<10$. Thus we may also take
$\eta_{j+1}=10x^2$. Finally $u_{j+1}\le6x^2\le10x^2$.
This proves \eqref{in:common} at the next level; the sequence remains
in $[0,.01]$. Solving its recurrence gives \eqref{in:xexplicit}.
\end{proof}
